\subsection{Reformulación transductiva}

Checkpoint~04 cambia la pregunta de modelado. Sea $\Train$ el conjunto etiquetado y
$\Target$ el conjunto fijo de 59 parcelas objetivo. La estimación final tiene forma
\[
\widehat{\bm y}_{\Target}
=
\mathcal A(X_{\Train},y_{\Train},X_{\Target}),
\]
donde $\mathcal A$ puede explotar la geometría conjunta de todos los $X$. Esto corresponde a
aprendizaje transductivo/semi-supervisado en el sentido de que el conjunto no etiquetado
específico participa en la representación \citep{Chapelle2006,Belkin2006}.

El primer paso no fue entrenar un modelo más complejo. Fue preguntar si la cercanía entre
parcelas realmente transfiere rendimiento y si los 59 objetivos forman un dominio distinto.

\subsection{Representaciones para topología}

Se construyeron tres espacios:
\begin{description}
\item[C0 base] Feature Table competition;
\item[C1 agronomic] las 350 features agronómicas;
\item[C4 all deterministic] base + agronomic + empirical, excluyendo expresiones
supervisadas descubiertas con $y$.
\end{description}

Para distancias multivariadas se imputa/estandariza usando exclusivamente $X$ y se obtiene un
embedding PCA de 20 componentes sobre las 197 parcelas. Esta operación es transductiva pero
$X$-only:
\[
Z=\mathrm{PCA}_{20}\bigl(\mathrm{scale}(X_{\Train}\cup X_{\Target})\bigr).
\]
No se usa $y$ para orientar los componentes.

\subsection{Clasificación adversarial train--target}

Se define una etiqueta auxiliar
\[
d_i=
\begin{cases}
0,&i\in\Train,\\
1,&i\in\Target.
\end{cases}
\]
Un clasificador cross-validated intenta distinguir ambos grupos a partir de C4/PCA. Si su AUC
estuviera muy por encima de 0.5, habría evidencia de covariate shift multivariado global. El
resultado fue
\[
\mathrm{AUC}=0.4892.
\]
Por tanto, los 59 objetivos no constituyen una población globalmente separable en la
representación examinada.

Este resultado no contradice que algunas parcelas individuales estén en bajo soporte. AUC
global cercano a azar implica solapamiento agregado, no que cada punto objetivo tenga un buen
vecino etiquetado.

\safefigure{../../reports/checkpoint_04a/figures/adversarial_roc.png}
{ROC de la clasificación adversarial train frente a target.}
{fig:cp04a-adversarial}

\subsection{¿Qué distancia transfiere rendimiento?}

Sobre las 138 etiquetas se construyeron todos los
\[
{138\choose2}=9{,}453
\]
pares. Para cada par $(i,j)$ se calcula
\[
\Delta y_{ij}=|y_i-y_j|
\]
y diversas métricas de disimilitud/similitud. La asociación se resume con Spearman, lo que
evita asumir linealidad entre distancia y diferencia de rendimiento.

\begin{table}[H]
\centering
\caption{Asociación entre similitud de covariables y diferencia absoluta de rendimiento.}
\label{tab:cp04a-transfer}
\begin{tabular}{lr}
\toprule
Métrica & Spearman con $|y_i-y_j|$ \\
\midrule
distancia geográfica & \textbf{0.4882} \\
distancia PCA C1 agronomic & 0.2200 \\
distancia PCA C4 deterministic & 0.1833 \\
distancia PCA C0 base & 0.1673 \\
DTW temporal & 0.0763 \\
Pearson temporal & -0.0503 \\
best-lag temporal correlation & -0.0482 \\
\bottomrule
\end{tabular}
\end{table}

Para distancias, una asociación positiva es la dirección útil: parcelas más lejanas tienden a
diferir más en rendimiento. Para similitudes tipo correlación, la dirección útil es negativa.
La geografía domina claramente a la similitud temporal mensual aislada.

\safefigure{../../reports/checkpoint_04a/figures/similarity_vs_yield_geo.png}
{Distancia geográfica frente a diferencia absoluta de rendimiento para pares etiquetados.}
{fig:cp04a-geo-transfer}

\subsection{Similitud temporal}

Se analizaron 11 trayectorias mensuales abril--octubre 2025: Sentinel-2 NDVI/EVI/LAI/FAPAR/
NDWI/MSI, Landsat VI6T y Planet NDVI/EVI/LAI/MSAVI. Para dos secuencias se calcularon
Pearson, Spearman, mejor correlación con lags de hasta dos meses y una distancia DTW
normalizada.

El resultado no es que la temporalidad carezca de información. Algunas parcelas tienen
análogos temporales convincentes. El hallazgo es más específico: \emph{como regla de
transferencia de rendimiento por sí sola}, la similitud mensual es mucho más débil que la
geografía. Por ello se evita construir el predictor final como nearest-neighbor temporal.

\subsection{Autocorrelación espacial de rendimiento}

La evidencia más fuerte de 04A proviene de Moran's $I$ \citep{Moran1950}. Para una matriz de
pesos espaciales $W=(w_{ij})$ y $z_i=y_i-\bar y$:
\begin{equation}
I
=
\frac{n}{\sum_{i,j}w_{ij}}
\frac{\sum_{i,j}w_{ij}z_i z_j}
{\sum_i z_i^2}.
\label{eq:moran}
\end{equation}

Sobre un grafo de cinco vecinos, el rendimiento observado produce
\[
I=0.6797,\qquad p_{\mathrm{perm}}=0.001,
\qquad \mathbb E_0[I]\approx-0.0073.
\]
Existe, por tanto, autocorrelación espacial positiva intensa.

\safefigure{../../reports/checkpoint_04a/figures/moran_yield.png}
{Moran scatter del rendimiento observado bajo el grafo espacial de cinco vecinos.}
{fig:cp04a-moran}

\subsection{Autocorrelación residual y dependencia del protocolo}

Los residuos OOF de los ensambles de 03C.3 muestran una diferencia reveladora:
\begin{table}[H]
\centering
\caption{Moran's $I$ de residuos OOF de la primera entrega.}
\label{tab:cp04a-residual-moran}
\begin{tabular}{lrr}
\toprule
Residuo & $I$ & $p_{\mathrm{perm}}$ \\
\midrule
E13 state-stratified & -0.0470 & 0.314 \\
E123 state-stratified & -0.0375 & 0.429 \\
E13 municipality-grouped & 0.5589 & 0.001 \\
E123 municipality-grouped & 0.5483 & 0.001 \\
\bottomrule
\end{tabular}
\end{table}

Cuando entrenamiento contiene parcelas cercanas, el modelo global absorbe parte de la
estructura espacial y sus residuos dejan de autocorrelacionarse fuertemente. Cuando se
retienen municipios, permanece una gran estructura espacial sin explicar. Esta observación
justifica modelos locales/grafo y explica la enorme brecha grouped de Checkpoint~03.

\subsection{Soporte de las 59 parcelas}

El diagnóstico 04A combina geografía, distancia C4, temporalidad, probabilidad adversarial y
consistencia de vecinos etiquetados en un score descriptivo. La taxonomía inicial fue:
\[
5\ \mathrm{A\ high}
+2\ \mathrm{B\ feature}
+5\ \mathrm{C\ geographic}
+13\ \mathrm{D\ extrapolation}
+34\ \mathrm{E\ mixed}.
\]

Entre los objetivos de mayor soporte aparecen AGC\_088, AGC\_190, AGC\_126, AGC\_128 y
AGC\_086. Entre los de menor soporte aparecen AGC\_173, AGC\_029, AGC\_120, AGC\_183 y
AGC\_003.

El score no es una probabilidad calibrada ni una predicción de rendimiento. Su función es
responder: ¿qué tan representado está un objetivo por la nube etiquetada?

\safefigure{../../reports/checkpoint_04a/figures/target_support_ranking.png}
{Ranking descriptivo de soporte de las 59 parcelas objetivo.}
{fig:cp04a-support}

\subsection{Decisión de transición}

Los resultados descartan dos rutas como prioridad: domain adaptation global agresiva y
nearest-neighbor temporal puro. En cambio, sostienen:
\begin{enumerate}
\item validación que reproduzca la geometría de las 59 parcelas;
\item local Ridge con vecindarios geográficos;
\item grafos sobre las 197 parcelas;
\item blends global/local sólo si reducen pseudo-test RMSE;
\item uso del soporte como hipótesis que debe validarse, no como regla manual.
\end{enumerate}

El siguiente checkpoint convierte estas hipótesis descriptivas en performance out-of-sample.
